\section{Introduction}
\label{sec:intro}

Multi-teacher on-policy distillation (MOPD) is used in LLM post-training to transfer specialized capabilities into a single model \citep{xiaomi2026mimov2flash,nvidia2026nemotron3ultra,kimiteam2026kimik3openfrontier,deepseekai2026deepseekv4highlyefficientmilliontoken}. The student generates responses, and a teacher assigned to each prompt domain supplies token-level supervision. The training objective is to improve several capabilities while preserving the student's existing performance. Achieving this requires deciding how much each domain and response contributes to the loss, and how many next-token alternatives the teacher supervises.

Recent work establishes both the potential of MOPD and the sensitivity of its results to these choices. \citet{ma2026mopdmultiteacheronpolicydistillation} study capability integration with policy-gradient and top-$k$ distillation objectives. Open-MOPD finds uneven transfer between domains and introduces domain token balancing, weighting by student--teacher performance gaps, and refreshed rewards \citep{gao2026openmopddiagnosingfixingcapability}. The Nemotron 3 Ultra report finds no advantage for top-$k$ or full-vocabulary matching over sampled-token supervision in its preliminary experiments \citep[Section~3.3.5]{nvidia2026nemotron3ultra}. Studies of OPD parameter changes also report concentration in a small set of stored weights \citep{yu2026denseupdates,shen2026opdgeometry}. These findings leave open how the loss weights and vocabulary targets affect the gradients, optimizer updates, and task scores of a jointly trained student.

Three distinctions matter for answering this question. First, equal numbers of prompts per domain do not imply equal domain contributions to a token-averaged loss, because response lengths differ. Even after fixing each domain's total loss weight, averaging over tokens or over responses gives different weights to long responses. Second, Adam updates depend on accumulated first- and second-moment estimates as well as the current gradient, and small FP32 updates may disappear when weights are stored in BF16. Sparse stored-weight changes or similar Adam update directions can therefore reflect the optimizer and numerical precision. Third, supervising more vocabulary tokens need not improve every task, even when the resulting parameter updates have similar concentration.

We study these distinctions through training runs and controlled gradient and optimizer comparisons with Qwen3-1.7B and four domain teachers. The main findings are:
\begin{enumerate}
\item \textbf{Different normalization rules produce distinct gradients but similar final top-64/Adam scores (Section~\ref{sec:normalization}).}
With equal prompt counts, mathematics receives 44.05\% of the token-averaged loss weight and instruction following 14.98\%. Fixing domain weights does not remove response-length weighting: token and response averaging have a mean gradient cosine of 0.680. Applying the same Adam moment estimates raises their update cosine to 0.960. The three normalization rules finish within 0.14 mean-score points under top-64 supervision, although their learning curves differ; sampled-token/SGD runs have a larger 1.07-point spread.
\item \textbf{BF16 sparsity and Adam update alignment can obscure the current teacher gradient (Section~\ref{sec:subnetworks}).}
Joint training changes about 97\% of FP32 master weights but only 7.3--8.5\% of stored BF16 weights. At the same student checkpoint, teacher-specific Adam updates have mean cosine similarity above 0.83 on their respective domain inputs. Removing the step that Adam would take using its accumulated moments with a zero current gradient reduces the cosine below 0.01. Thus aligned optimizer steps do not establish aligned current teacher gradients.
\item \textbf{More vocabulary supervision does not consistently improve task performance (Section~\ref{sec:density}).}
Sampled-token, top-64, and full-vocabulary losses produce similarly concentrated single-step updates when the top-64 intersection already contains over 99.9\% of student probability on average. In joint training, top-64 supervision scores 2.60 points higher on mathematics and 1.77 points lower on GPQA than sampled-token supervision at update 500; all paired task intervals include zero. Its mean advantage also shrinks from 1.69 points at update 250 to 0.29 points at update 500.
\end{enumerate}

These results motivate evaluating supervision choices with per-task learning curves in addition to gradient and parameter statistics. Section~\ref{sec:recipe} evaluates development-selected checkpoints on separate test sets: adding a per-domain retention constraint does not change any selected checkpoint in the ten tested trajectories.
